\clearpage
\VerdierPhysicalPageAnchor{0117}%
\begingroup
\small
\setlength{\abovedisplayskip}{4pt plus 2pt minus 2pt}%
\setlength{\belowdisplayskip}{4pt plus 2pt minus 2pt}%

\VerdierSemanticAnchor{ch2-1-2-5-combined-diagram}
したがって、次の図式は可換である。
\begin{equation*}
  \resizebox{0.98\linewidth}{!}{%
    $\begin{tikzcd}[ampersand replacement=\&,column sep=2.1em,row sep=3.1em]
      X\oplus X' \arrow[r,"{u\oplus u'}"] \arrow[d,"{\operatorname{id}}"']
        \& Y\oplus Y' \arrow[r,"{v\oplus v'}"] \arrow[d,"{\operatorname{id}}"']
        \& Z\oplus Z' \arrow[r,"{w\oplus w'}"] \arrow[d,"{(f,f')}"']
        \& (X\oplus X')[1] \arrow[d,"{\operatorname{id}}"] \\
      X\oplus X' \arrow[r,"{u\oplus u'}"']
        \& Y\oplus Y' \arrow[r,"m"']
        \& L \arrow[r,"n"']
        \& (X\oplus X')[1]
    \end{tikzcd}$%
  }\,.
\end{equation*}

\VerdierSemanticAnchor{ch2-1-2-5-sum-cone-map}
系を証明するには、射
\[
  (f,f')\colon Z\oplus Z'\longrightarrow L
\]
が同型であることを示せばよい。したがって、任意の対象 $M$ に対して射

\VerdierSemanticAnchor{ch2-1-2-5-hom-test}
\[
  \operatorname{Hom}_{\mathcal D}\bigl(M,(f,f')\bigr)
\]
がアーベル群の同型であることを示せば十分である。これは (1.2.1) と
五項補題から直ちに従う。

\VerdierSemanticAnchor{ch2-1-2-5-proof-complete}

\VerdierSemanticAnchor{ch2-1-2-6}
\medskip
\noindent\textbf{系 1.2.6.}\quad
$X\xrightarrow{u}Y\xrightarrow{v}Z\xrightarrow{w}X[1]$ を
$w=0$ を満たす $\mathcal D$ の完全三角とする。このとき、射 $v$ の切断
$s\colon Z\longrightarrow Y$ が存在する。任意の切断
$s\colon Z\longrightarrow Y$ は同型

\VerdierSemanticAnchor{ch2-1-2-6-split-isomorphism}
\[
  X\oplus Z\xrightarrow[\sim]{(u,s)}Y\,.
\]

\VerdierSemanticAnchor{ch2-1-2-6-converse}
逆に、次の形の三角

\VerdierSemanticAnchor{ch2-1-2-6-1}
\begin{equation}
  X\xrightarrow{\binom{\operatorname{id}}{0}}X\oplus Z
  \xrightarrow{(0,\operatorname{id})}Z\xrightarrow{0}X[1]
  \tag{1.2.6.1}
\end{equation}
はすべて完全三角である。

\VerdierSemanticAnchor{ch2-1-2-6-third-assertion}
第三の主張は TRI、TRII、および (1.2.5) から直ちに従う。

\VerdierSemanticAnchor{ch2-1-2-6-first-assertion}
第一の主張を証明しよう。(1.2.1) により、任意の $W$ に対して列

\VerdierSemanticAnchor{ch2-1-2-6-2}
\begin{equation}
  0\longrightarrow\operatorname{Hom}_{\mathcal D}(W,X)
  \longrightarrow\operatorname{Hom}_{\mathcal D}(W,Y)
  \longrightarrow\operatorname{Hom}_{\mathcal D}(W,Z)
  \longrightarrow0
  \tag{1.2.6.2}
\end{equation}
\endgroup
